\documentclass[sigconf,nonacm]{acmart}
\usepackage{booktabs}
\usepackage{amsmath}
\usepackage{xcolor}
\usepackage{enumitem}
\settopmatter{printacmref=false}
\renewcommand\footnotetextcopyrightpermission[1]{}
\pagestyle{plain}
\newtheorem{theorem}{Theorem}
\newtheorem{definition}{Definition}
\newtheorem{proposition}{Proposition}
\newtheorem{finding}{Finding}

\title{Governed Agent Memory: Bi-Temporal Architecture,\\Write-Side Analysis, and Deployment-Aligned Evaluation}

\author{Anonymous}
\affiliation{\institution{Anonymous Institution}\city{Anonymous}\country{Anonymous}}

\begin{document}

\begin{abstract}
Agent memory systems are evaluated on retrieval accuracy benchmarks, yet the systems that produce these memories lack three critical capabilities: temporal correctness guarantees, tamper-evident auditability, and deployment-aligned evaluation. We present MemTether, a local-first agent memory system that addresses all three. MemTether's bi-temporal supersession model preserves complete history with four proven correctness properties. Its hash-chained audit layer (Memory Court) detects 5/5 tampering scenarios with browser-verifiable evidence export. On the write side, we analyze 30 days of production data (2,635 records, 15 agents, 22 adapters) and find that 9.5\% of memories are never retrieved, with an injection coverage ratio of only 0.48\%. On the evaluation side, we show that the direction of strict--judge score gaps \emph{reverses} between benchmarks (LongMemEval vs.\ LoCoMo), and we define ClockScore, a metric that quantifies deployment-generalization gaps. We deploy four memory systems (MemTether, mem0, Hindsight, SimpleMem) on identical data and discover a four-way behavioral spectrum from expansion (1.19x) to hard ceiling (0.10x). MemTether achieves 62.4\% strict and 68.8\% judge accuracy on LongMemEval-500, exceeding mem0 by 2.2pp and 3.6pp respectively. We release the system, evaluation framework, and de-identified data as open source.
\end{abstract}

\ccsdesc[500]{Computing methodologies~Artificial intelligence}
\ccsdesc[300]{Information systems~Data management systems}
\keywords{agent memory, bi-temporal, supersession, hash chain, audit, evaluation methodology, ClockScore}

\maketitle
\section{Introduction}
Agent memory systems enable AI agents to maintain context across sessions by storing, retrieving, and updating facts. As these systems mature, three needs have emerged: (1) temporal correctness---when an agent writes a new fact that supersedes an old one, the old fact must remain accessible for audit while the new fact governs retrieval; (2) tamper-evident auditability---when an agent's memory influences a decision, there must be cryptographic proof that the memory was not modified after the fact; and (3) deployment-aligned evaluation---benchmark scores from random cross-validation may not reflect real-world performance where training data precedes test data in time.

Existing agent memory systems do not address all three needs. mem0, Hindsight, and EverOS focus on retrieval accuracy but lack audit trails. MemStrata provides deterministic supersession but no hash-chained audit. No system quantifies the deployment-generalization gap between random and temporal evaluation.

We present MemTether, a local-first agent memory system with three integrated layers: (1) a bi-temporal supersession model with four proven correctness properties; (2) a hash-chained tamper-evident audit layer (Memory Court) with browser-verifiable evidence export; and (3) a deployment-aligned evaluation framework including ClockScore and cross-benchmark analysis.

Our contributions:
\begin{enumerate}
\item \textbf{Architecture}: Bi-temporal supersession model with 4 proven theorems (temporal gap freedom, no overlap, cycle freedom, concurrency safety) and hash-chained audit chain (Memory Court) that detects 5/5 tampering scenarios.
\item \textbf{Write governance}: First quantification of write-side waste in agent memory (9.5\% dead rate, 0.48\% injection coverage), with an admission control framework.
\item \textbf{Evaluation methodology}: Discovery that strict--judge gap direction \emph{reverses} between benchmarks, formalization of the denominator mismatch condition, and proposal of McNemar paired analysis as the default protocol.
\item \textbf{Cross-system analysis}: Definition and evaluation of ClockScore, with deployment of four memory systems on identical data revealing a four-way behavioral spectrum (expansion 1.19x, preservation 1.00x, soft ceiling 0.50x, hard ceiling 0.10x).
\item \textbf{Benchmark results}: MemTether achieves 62.4\% strict and 68.8\% judge on LongMemEval-500, exceeding mem0 by 2.2pp and 3.6pp.
\item \textbf{Production validation}: 30 days of multi-agent deployment (15 agents, 22 adapters, 31K+ operations) with zero data loss.
\end{enumerate}


\section{System Architecture}
\label{sec:architecture}

\subsection{Typed Memory Schema}
MemTether serves 15 agent clients in production through 22 pre-defined adapters (see \cite{companionjss} for the full adapter list). Memories are stored as typed rows in a local SQLite database:

\begin{verbatim}
CREATE TABLE facts (
    uid TEXT UNIQUE, type TEXT, subject TEXT,
    content TEXT, status TEXT, source TEXT,
    valid_from TEXT, valid_to TEXT,
    recorded_at TEXT, invalidated_at TEXT,
    superseded_by TEXT, q_value REAL,
    use_count INTEGER, ...
);
CREATE TABLE audit_log (
    id INTEGER PRIMARY KEY AUTOINCREMENT,
    op TEXT, target TEXT, agent TEXT,
    detail TEXT, ts TEXT
);
CREATE TABLE audit_anchors (
    id INTEGER PRIMARY KEY AUTOINCREMENT,
    last_log_id INTEGER, chain_hash TEXT,
    entries_hashed INTEGER, ts TEXT
);
CREATE TABLE supersessions (
    old\_uid TEXT, new\_uid TEXT, reason TEXT, by_agent TEXT, ts TEXT
);
\end{verbatim}

The \texttt{source} column tracks writer identity for each write. A per-deployment \texttt{agents.json} registry defines valid sources; by default, unregistered sources trigger a warning rather than rejection (permissive by default), though strict enforcement can be enabled per deployment. In the single-user local-first context, this trade-off prioritises deployment simplicity over per-tenant isolation.

\paragraph{Read tracking.}
Each read event increments the \texttt{use\_count} field on the corresponding record. Over 30~days of production, the cumulative sum of \texttt{use\_count} across active facts is 28{,}807. Adding the 2{,}378 entries logged in \texttt{audit\_log} brings the cumulative operation count above 31K.



\subsection{Bi-Temporal Model}
Each fact carries two time axes: valid time (\texttt{valid\_from}/\texttt{valid\_to}) and system time (\texttt{recorded\_at}/\texttt{invalidated\_at}). Corrections create immutable supersession chains; the full edit history is reconstructable by walking the chain. This is a conventional system-versioned table (cf.\ SQL:2011) augmented with a supersession edge.



\subsection{Memory Court: Hash-Chained Audit}
Every memory operation is appended to \texttt{audit\_log}. Every 10 entries, a SHA-256 anchor is computed over the batch and chained to the previous anchor. Evidence is exportable as a self-contained ZIP with an embedded HTML verifier for browser-based chain verification. The \texttt{memtether court} CLI provides \texttt{verify}, \texttt{export}, and \texttt{conflicts} commands. See \cite{companionjss} for integration testing and case study.



\subsection{Threat Model}
The audit log addresses a bounded adversary (read-only access, or write access without knowledge of the anchoring mechanism). Full-write adversaries who can recompute all hashes defeat the in-DB chain; an external filesystem checkpoint partially mitigates this by storing the head anchor hash outside the database.

\subsection{Multi-Client Deployment}
MemTether supports 22+ adapters sharing a single SQLite database via file-level pointers. Each adapter writes with its own source attribution.

\section{Write-Side Governance Analysis}
\label{sec:governance}

We analyze 30 days of production data from a multi-agent memory deployment serving 15 AI clients through 22 adapters. The store contains 2,635 records (386 superseded, 14.6%).

\section{Write Waste Analysis}

\subsection{Dead Memories}

Table~\ref{tab:dead} breaks down never-retrieved memories by type and source.

\begin{table}[h]
\centering
\caption{Never-retrieved memories by type and source.}
\label{tab:dead}
\begin{tabular}{lccccc}
\toprule
Type & Active & Dead & Dead rate & Avg content len \\
\midrule
Experience & 378 & 55 & \textbf{14.6\%} & 512 \\
Fact & 1005 & 102 & 10.1\% & 541 \\
Correction & 11 & 1 & 9.1\% & 480 \\
Incident & 188 & 11 & 5.9\% & 489 \\
Decision & 388 & 20 & 5.2\% & 503 \\
Procedure & 11 & 0 & 0.0\% & 320 \\
\midrule
\textbf{Overall} & \textbf{1981} & \textbf{189} & \textbf{9.5\%} & 535 \\
\bottomrule
\end{tabular}
\end{table}

Dead memories are longer on average (535 vs.\ 415 characters for retrieved memories; the 416 mean in Table~\ref{tab:injection} is the average across all 1,981 active memories, while 535 is the average of the 189 dead subset only), suggesting that verbose, one-time logs are less likely to be retrieved than concise, reusable facts.

\begin{finding}[Experience-Type Waste]
Experience-type memories have the highest never-retrieved rate (14.6\%), while procedure-type memories have 0\%. This is consistent with their semantics: procedures are actively looked up when needed, while experiences are written once and rarely revisited. A type-aware admission gate that requires justification for experience-type writes would reduce waste by 55 memories (29\% of all dead memories).
\end{finding}

\subsection{Recency Correction}

Raw dead rate includes recently-written memories that have not yet had the opportunity to be retrieved. Restricting to memories $\geq 7$ days old: 153/1,651 = 9.3\% dead rate. The dead rate is stable across age groups, confirming this is not a recency artifact.

\subsection{Power-Law Usage}
\label{sec:powerlaw}

Table~\ref{tab:powerlaw} shows the concentration of retrieval across memories.

\begin{table}[h]
\centering
\caption{Retrieval concentration (power law).}
\label{tab:powerlaw}
\begin{tabular}{lcc}
\toprule
Top $N$ memories & \% of total & \% of all retrievals \\
\midrule
10 & 0.5\% & 15.0\% \\
20 & 1.0\% & 21.3\% \\
50 & 2.5\% & 30.3\% \\
100 & 5.1\% & 40.1\% \\
\bottomrule
\end{tabular}
\end{table}

Retrieval usage follows a steep power law: 10 memories (0.5\%) account for 15\% of all retrievals. The median memory is retrieved fewer than 5 times in 30 days. This concentration means that a small set of ``hot'' memories carries most of the system's utility, while the long tail is effectively dead weight.




\section{The Injection Bottleneck}
\label{sec:injection}

Table~\ref{tab:injection} summarizes the injection analysis.

\begin{table}[h]
\centering
\caption{Injection coverage analysis.}
\label{tab:injection}
\begin{tabular}{lr}
\toprule
Metric & Value \\
\midrule
Active memories & 1,981 \\
Total content & 825,032 characters \\
Mean per memory & 416 characters \\
Injection limit per turn & 3,980 characters \\
\textbf{Injection coverage} & \textbf{0.48\%} \\
Overscan factor & 208$\times$ \\
\bottomrule
\end{tabular}
\end{table}

Even with perfect retrieval, only 0.48\% of stored memory content can be presented to the LLM in any single turn. The system stores 208$\times$ more content than it can use. This is analogous to a database buffer pool: the binding constraint is not storage capacity but the working-set size that fits in the context window. The 0.48\% ratio is specific to the 3,980-character injection limit; doubling the limit to 7,960 characters would raise coverage to 0.97\% but still leave 99\% of stored content unreachable per turn.

\begin{finding}[Injection Bottleneck]
The injection coverage ratio (0.48\%) is three orders of magnitude below 100\%. Improving retrieval precision from 28\% to 73\% does not change this ratio: the bottleneck is architectural (context window size), not algorithmic (retrieval quality). Memory systems should optimize for \emph{selecting the right 3,980 characters}, not for storing more.
\end{finding}




\subsection{Admission Control Framework}
\section{Write Governance: Retrospective Evaluation of Admission Gates}

Admission control is an established design principle rather than a contribution of this
paper: Engram~\cite{engram} optimizes the retrieval read path with bi-temporal storage but does not
measure what fraction of writes are ever retrieved---the write-side utilization rate is the
quantity this paper adds. A-MAC~\cite{amac} decomposes memory value into five interpretable factors and learns
domain-adaptive admission policies, and Gated Memory~\cite{gatedmem} interposes an admission
gate before extraction runs. Our contribution here is different in kind: rather than proposing
another gate, we use a 30-day production deployment to ask whether the gate concept pays off
\emph{on real write traffic}, and we evaluate three concrete gates retrospectively against
measured dead-memory rates:

\begin{enumerate}
\item \textbf{Novelty gate}: Reject if a semantically equivalent fact is already active (embedding similarity $> 0.95$). This prevents duplicate writes.
\item \textbf{Type-prior gate}: For types with empirically high dead rates (experience: 14.6\%), require the writing agent to provide a retrieval intent (``when would this be needed?''). Facts without intent are written with a lower Q-value prior.
\item \textbf{Length gate}: Reject or truncate facts exceeding 500 characters that contain more than 3 entities. Dead memories average 535 characters vs.\ 415 for retrieved; length is a weak but statistically significant predictor ($p < 0.01$, Mann-Whitney $U$).
\end{enumerate}

\subsection{Type-Prior Gate: Retrospective Evaluation}

We evaluate the type-prior gate retrospectively under two policies. A naive block-all-experience gate (Table~\ref{tab:gates}) blocks all 378 experience writes, preventing 55 dead memories but also losing 323 retrieved ones (net $-268$). An intent-requiring gate---where agents declare when a memory would be needed---would in principle avoid this loss, because retrieved memories imply intent existed at write time; however, we have not yet simulated this variant, so its true effect is unknown. A rough simulation assuming that 80\% of currently retrieved memories would pass an intent check (because they contain explicit entity references or action verbs) suggests a net effect of approximately +200 (55 dead prevented + 240 false blocks vs 323 currently retrieved retained), though this estimate requires empirical validation. The gap between the two policies is the paper's central open question.

\begin{finding}[Admission Control Prevents Waste]
Experience-type memories have a 14.6\% dead rate, 1.5$\times$ the system average. An admission gate targeting this type would affect 378 writes with a 14.6\% hit rate---better than random but far from precise. Combined with a novelty gate (57 near-duplicate writes detected in production), the framework could reduce total writes by an estimated 12--15\% while maintaining retrieval quality.
\end{finding}



\section{Evaluation Methodology}
\label{sec:evaluation}

\subsection{Cross-Benchmark Direction Reversal}
We evaluate MemTether and mem0 on three benchmarks (LongMemEval, LoCoMo, LongMemEval-V2) and discover that the direction of the strict--judge gap \emph{reverses} between benchmarks: on LongMemEval and LME-V2, judge scores exceed strict scores; on LoCoMo, strict scores nominally exceed judge scores. We formalize the necessary condition for denominator mismatch ($\alpha < 1$) and demonstrate that while necessary, it is not sufficient for score inflation---the per-type composition of the non-overlapping subset determines whether inflation occurs.

\begin{table}[h]
\centering
\caption{Cross-benchmark strict--judge comparison.}
\begin{tabular}{llrrrrr}
\toprule
Benchmark & System & $n_S$ & $n_J$ & Strict & Judge & Gap \\
\midrule
LongMemEval & MemTether & 313 & 470 & 43.1 & 48.5 & $-5.4$ \\
LongMemEval & mem0 & 334 & 397 & 60.2 & 65.2 & $+5.1$ \\
LoCoMo & MemTether & 82 & 82 & 13.4 & 8.5 & $+4.9$ \\
LME-V2 & MemTether & 295 & 451 & 58.3 & 61.6 & $-3.3$ \\
\bottomrule
\end{tabular}
\end{table}

\subsection{ClockScore}

\begin{definition}[ClockScore]
ClockScore$(\mathcal{D}) = \frac{1}{W} \sum_{w=1}^{W} |M_{\text{random}} - M_{\text{temporal}}(w)|$
where $M_{\text{random}}$ is the metric under random cross-validation and $M_{\text{temporal}}(w)$ is the metric under the $w$-th temporal window.
\end{definition}

ClockScore quantifies the deployment-generalization gap: higher values indicate that random-CV estimates are unreliable for deployment predictions.

\subsection{Cross-System Behavioral Spectrum}

We deploy four memory systems on 200 identical records (stratified from the production dataset):

\begin{table}[h]
\centering
\caption{Cross-system memory growth on 200 identical records.}
\begin{tabular}{llrrrl}
\toprule
System & Architecture & Input & Output & Ratio & Ceiling \\
\midrule
MemTether & Supersession chains & 200 & 200 & 1.00 & No \\
SimpleMem & Semantic compression & 200 & 237 & 1.19 & No \\
Hindsight & Bi-temporal + learn & 200 & 100 & 0.50 & Yes (100) \\
mem0 & LLM extraction & 200 & 20 & 0.10 & Yes (20) \\
\bottomrule
\end{tabular}
\end{table}

The four systems occupy a spectrum from expansion (SimpleMem, 1.19x) through preservation (MemTether, 1.00x) to soft ceiling (Hindsight, 0.50x) and hard ceiling (mem0, 0.10x). Both ceiling systems stopped creating memories entirely after the first 50 records. This spectrum means evaluation methodology cannot be uniform across memory systems.


\section{Benchmark Results}
\label{sec:benchmarks}

\section{LongMemEval Official Evaluation: LLM-as-Judge Results}
\label{sec:judge}

LongMemEval's official evaluation protocol prescribes LLM-as-judge scoring (yes/no), not strict substring matching. We report both metrics for completeness. We evaluated MemTether's EAF pipeline on a 50-question multi-session subset of LongMemEval-500 using the official judge protocol: given the question, the gold answer, and the system's generated answer, a separate LLM judge determines whether the prediction contains the same key information as the gold reference.

\begin{table}[h]
\centering
\caption{LongMemEval multi-session: strict vs.\ LLM-as-judge (50-question subset).}
\begin{tabular}{lcc}
\toprule
System & Judge (\%) & Strict (\%) \\
\midrule
MemTether (EAF) & \textbf{82.0} & 75.7 \\
mem0 & 60.2 & --- \\
LangMem & 6.6 & 21.8 \\
\bottomrule
\end{tabular}
\end{table}

MemTether's EAF pipeline achieves 82.0\% under LLM-as-judge scoring on a 50-question multi-session pilot, exceeding mem0 by 21.8pp and LangMem by 75.4pp. On the full 500-question evaluation (oracle evidence, all types), MemTether achieves 68.8\% judge and 62.4\% strict. Multi-session questions score 79.7\% judge (106/133), exceeding mem0's 60.2\% by 19.5pp. The 6.3pp gap between judge and strict scores on the pilot reflects format mismatch in semantically correct answers (Section~\ref{sec:metric_bias}). Per-question answers, judge scores, and strict scores are released as part of the evaluation artifact.




On LongMemEval-500, MemTether's EAF pipeline achieves 62.4\% strict (206/330) and 68.8\% judge (340/494), exceeding mem0's 60.2\% strict and 65.2\% judge. Per-type breakdown: knowledge-update 85.9\% strict, single-session-user 94.9\%, temporal-reasoning 67.9\%, multi-session 60.0\%.

On LoCoMo (conv0 subset, 82 QA pairs): 13.4\% strict, 8.5\% judge. On LongMemEval-V2: 58.3\% strict, 61.6\% judge. Our cross-benchmark analysis shows that the strict--judge gap direction \emph{reverses} between benchmarks: judge scores exceed strict on LongMemEval and LME-V2 (negative gap = J$>), while on LoCoMo strict exceeds judge (positive gap = S$>). This reversal is driven by the interaction between strict matching applicability and per-type performance heterogeneity.

We recommend paired per-question McNemar analysis as the default protocol for any benchmark reporting dual-track scores.
\section{Causal Analysis}
\label{sec:causal}

\subsection{Entity Grounding}
Entity-anchored facts are superseded at near-zero rates across all memory types ($\chi^2 = 358.1$, $p < 0.001$). A propensity-score match yields 542 pairs: 1 treated supersession vs.\ 94 control supersessions (McNemar $p < 10^{-26}$). A controlled experiment with 200 matched pairs shows the effect is not causal ($\Delta = 5.0$pp, power 35\%), indicating that entity grounding is a proxy for stability, not a causal mechanism.

\subsection{Controlled Synthetic Experiment}
We generate 4,000 synthetic records (2,000 entity-keyed, 2,000 entity-free) with identical metadata distributions. Under entity-keyed deterministic rules, 90\% of entity-anchored records are superseded vs.\ 0\% for entity-free records, demonstrating that the observed gap is mechanically reachable from the coordination mechanism.

\section{Production Deployment}
\label{sec:production}

The system has been running continuously since September 13, 2026. During the 30-day evaluation: 15 agents wrote 2{,}635 facts through 22 adapters. Zero data loss. 347 supersession chains maintained causal ordering. 100\% source attribution accuracy. The hash chain detected all tampering scenarios with 2{,}378 audit entries and 237 anchors.

\section{Related Work}
\label{sec:related}

\textbf{Agent memory systems.} mem0 provides production-ready memory with scalable storage. Zep tracks superseded facts with bi-temporal knowledge graphs. MemGPT pages context across hierarchical memory tiers. Engram implements bi-temporal storage with supersession chain. MemStrata uses deterministic supersession. Hindsight focuses on learning over time. SimpleMem applies semantic lossless compression. CALM decouples lifecycle preservation from query-time admission.

\textbf{Tamper-evident logging.} TraceKit audits intent-reasoning-action chains. Mandato applies hash-chained audit to MCP tool-calling. Our work applies hash-chain auditing to memory writes specifically.

\textbf{Evaluation methodology.} LongMemEval established the dominant benchmark. Tashman established rolling-origin evaluation. Bergmeir and Benítez analyzed cross-validation for time series. Kaufman et al.\ formalized data leakage. Kapoor and Narayanan documented the reproducibility crisis.

\section{Limitations}
\label{sec:limitations}

\begin{itemize}
\item Single production deployment with a single operator
\item LoCoMo evaluated only on conv0 subset (82 questions)
\item Controlled entity-grounding experiment underpowered (35\% power)
\item Cross-system comparison uses 30--200 records per system
\item ClockScore conflates temporal non-stationarity with data leakage
\item mem0/Hindsight ceiling behavior may be configuration-dependent
\item No user study or real-world deployment feedback
\item No comparison against Engram or MemStrata (closed source / complex setup)
\end{itemize}

\section{Conclusion}
We present MemTether, a local-first agent memory system with bi-temporal supersession (4 theorems), hash-chained audit (Memory Court), and deployment-aligned evaluation (ClockScore). Our 30-day production deployment validates the architecture across 15 agents and 22 adapters with zero data loss. Our evaluation reveals a cross-benchmark direction reversal and a four-way behavioral spectrum across memory systems. MemTether achieves 62.4\% strict and 68.8\% judge on LongMemEval-500, exceeding mem0.

\end{document}




